\chapter{Hyphema (Bleeding in the Eye)}

\section{Definition}
Hyphema is the accumulation of blood in the anterior chamber of the eye between the cornea and the iris. It most commonly results from blunt trauma but can occur spontaneously in certain conditions. Hyphema is an ocular emergency.

\section{Pathophysiology}
Blunt trauma to the eye causes a sudden increase in intraocular pressure that tears the blood vessels of the ciliary body or iris root. Blood leaks into the anterior chamber, where it settles by gravity. The blood can obstruct the trabecular meshwork, impairing aqueous humor drainage and increasing intraocular pressure. Rebleeding can occur 3-5 days after the initial injury.

\section{Etiology}
\begin{itemize}
\item Blunt eye trauma (most common cause, especially in children)
\item Penetrating eye injury
\item Sports-related eye injuries (racquetball, baseball, paintball)
\item Spontaneous hyphema: iris neovascularization (diabetes), intraocular tumors, bleeding disorders
\item Anticoagulant or antiplatelet therapy (warfarin, aspirin, clopidogrel)
\item Sickle cell disease (increases risk of elevated IOP)
\item Post-surgical (cataract surgery, glaucoma surgery)
\item Child abuse (shaken baby syndrome)
\end{itemize}

\section{Indications for Evaluation}
\begin{itemize}
\item Visible blood settling in the lower part of the eye (anterior chamber)
\item Blurred or decreased vision
\item Eye pain
\item Sensitivity to light (photophobia)
\item Elevated intraocular pressure
\item Blood may layer (microhyphema) or fill the entire anterior chamber (total hyphema, eight-ball hyphema)
\item History of trauma or bleeding risk factors
\end{itemize}

\section{Related Symptoms}
\begin{itemize}
\item Eye pain
\item Blurred vision
\item Traumatic iritis
\item Glaucoma (elevated intraocular pressure)
\item Corneal blood staining (long-term complication)
\item Orbital fracture (if associated with trauma)
\end{itemize}

\section{Self-Care/Home Management}
\begin{itemize}
\item Seek emergency eye care immediately
\item Keep the head elevated 30-45 degrees
\item Wear an eye shield or patch
\item Do NOT rub the eye
\item Avoid strenuous activity and heavy lifting
\item Do not take aspirin or NSAIDs (may increase bleeding)
\item Use a protective shield during sleep
\end{itemize}

\section{Diagnostic Approach}
\begin{itemize}
\item Slit-lamp examination to visualize blood in the anterior chamber
\item Grading of hyphema (grade 0 to grade 4 based on blood level)
\item Intraocular pressure measurement (tonometry)
\item Detailed eye examination with dilation (after ruling out globe rupture)
\item CT scan if orbital fracture or intraocular foreign body suspected
\item Blood coagulation profile if spontaneous or recurrent
\end{itemize}

\section{Treatment/Management Overview}
\begin{itemize}
\item Elevate head of bed 30-45 degrees to promote settling
\item Topical corticosteroids to reduce inflammation
\item Cycloplegics (atropine) for comfort and to prevent synechiae
\item Anti-glaucoma medications for elevated IOP
\item Aminocaproic acid (tranexamic acid) to prevent rebleeding
\item Anterior chamber washout if IOP is uncontrolled or total hyphema
\item Surgical evacuation for persistent clot or corneal blood staining
\item Treat underlying cause (sickle cell, coagulation disorder)
\end{itemize}
